\chapter{Introduction}
\label{chap:introduction}

\section{Motivation}
\label{sec:introduction-motivation}

Sorting normally starts with an arbitrary sequence of keys. Jordan sorting
starts with more information: the keys are coordinates of intersections
encountered while traversing a simple plane curve. Their input order along the
curve and their desired order along a line differ, but the absence of
self-intersections restricts how consecutive curve segments can be nested.
That geometric restriction permits specialized recognition and sorting
algorithms with a theoretical linear-time bound
\parencites{hoffmann1986sorting}{fung1990simplified}.

The published bound, however, combines an algorithmic control flow with
specialized dynamic-list representations. The 1986 algorithm uses
level-linked search trees, while the 1990 simplification relies on
heterogeneous finger trees to support local list operations efficiently. A
direct implementation with ordinary arrays or lists can preserve the
high-level algorithm while having different operation costs. Moreover, the
1990 presentation is concise: translating its incremental steps into
executable code requires explicit choices about pair orientation, sibling-list
ownership, geometric endpoints, and the exceptional role of the first point.

This thesis studies that implementation gap. It reconstructs the 1990
Step~1/2/3 sorting procedure as an executable Python core, maintains its own
partial sorted order and family state, and evaluates the resulting
ordinary-list implementation on controlled, oracle-certified valid inputs.
The aim is not to reproduce the historical linear-time data structure or to
infer asymptotic complexity from runtime measurements. It is to make the
paper-facing control flow concrete, auditable, and empirically observable
while keeping the boundary between theorem, implementation, and experiment
explicit.

\section{Problem and Scope}
\label{sec:introduction-scope}

For a Jordan sequence $z_1,\ldots,z_n$, the computational task considered here
is to recover increasing coordinate order from curve-traversal order. The
geometric model and its nested upper and lower pair families are introduced in
\cref{chap:background}; the executable reconstruction is specified in
\cref{chap:algorithm}.

The project contains two deliberately different research implementations. The
\emph{simplified reference pipeline} performs oracle validation, constructs
structural information, and returns oracle-derived sorted output. The
\emph{paper ordinary-list core} instead accepts a pre-certified valid input,
executes the reconstructed incremental procedure, and recovers its output from
the maintained partial order. It does not consume the oracle's sorted result.
An outer wrapper may certify the input before calling the core, and the
experiment may compare its return value with the oracle afterward; both
operations are outside the timed paper call. This provenance boundary is
formalized in \cref{sec:implementation-certification}.

The empirical scope is valid-input sorting. Recognition of arbitrary or
invalid sequences, error correction, polygon clipping, and the historical
theoretical backends are not evaluated. Sibling lists use ordinary Python
lists, so scans, slicing, copying, and ownership rebinding remain part of the
implemented cost. Consequently, the historical linear-time theorem does not
transfer to this implementation.

\section{Research Questions}
\label{sec:introduction-questions}

Within those boundaries, the thesis addresses three questions:

\begin{enumerate}[label=\textbf{RQ\arabic*:},leftmargin=*]
  \item \textbf{Executable reconstruction and correctness.} Can an
    ordinary-list reconstruction of the 1990 incremental sorting procedure
    independently recover sorted order for oracle-certified valid inputs while
    maintaining the required partial-order and sibling-family invariants?

  \item \textbf{Observed runtime behavior.} Across the frozen input sizes and
    controlled families, how does the pre-certified paper call in
    \texttt{minimal} mode behave relative to Python sort and the complete
    oracle-backed reference pipeline under their explicitly different timing
    scopes?

  \item \textbf{Structure and implementation cost.} Within the controlled
    sample, what descriptive relationships appear between structural
    properties, checked operation counters, and the runtime of the ordinary-list
    implementation?
\end{enumerate}

The first question is addressed through focused tests, exhaustive repository
validation over bounded oracle-certified cases, checked state audits, and
deterministic replay. The second uses the frozen Week~12 formal experiment.
The third is exploratory: its associations are descriptive rather than causal
or asymptotic.

\section{Contributions}
\label{sec:introduction-contributions}

The work presented in this thesis makes four implementation and empirical
contributions.

\begin{enumerate}
  \item It gives an executable reconstruction of the 1990 Step~1/2/3 control
    flow, informed by the 1986 algorithm where the later presentation leaves
    operational details implicit. The reconstruction records its endpoint,
    orientation, and first-point clarifications rather than presenting them as
    the only possible reading of the paper.

  \item It implements separate partial-order and sibling-family structures
    with explicit ownership, transactional splitting, rollback, local safety
    checks, and full invariant auditing. The paper core recovers output from
    this maintained state rather than from the oracle-backed reference path.

  \item It separates optional observation and complete audit machinery from
    the timed path through fixed execution policies. Production sorting and
    checked diagnostics share one core control flow, while deterministic replay
    and complete state validation provide an untimed audit of recorded state.
    The \texttt{minimal} policy still retains local safety, stage state, and
    final output recovery.

  \item It establishes a reproducible empirical workflow with frozen cases,
    documented timing scopes, immutable run evidence, and an independent
    fail-closed validator. The resulting study compares Python sort, the
    oracle-backed reference pipeline, and the paper ordinary-list call without
    treating their scopes as interchangeable.
\end{enumerate}

These contributions concern reconstruction, correctness engineering, and
empirical characterization. They do not include level-linked search trees,
heterogeneous finger trees, a theoretical linear-time split/update backend, or
a proof of linear running time for the Python implementation.

\section{Headline Findings}
\label{sec:introduction-findings}

All 3600 measured rows in the frozen Week~12 run have oracle-certified inputs,
correct output, passed checked audits, and no recorded error. The 180
case-algorithm summaries and all 60 case audits also pass. This is empirical
correctness evidence for the frozen cases, not a proof for every Jordan
sequence.

The median exact-case paper/reference runtime ratio decreases across the five
tested sizes, from $3.226$ at $n=32$ to $0.567$ at $n=512$. It is above one
through $n=128$ and below one at $n=256$ and $n=512$. This ratio compares two
different timed pipelines: the reference value includes its complete
oracle-backed reference work, whereas the paper value includes only the
pre-certified \texttt{minimal} sorting call, with certification and checked
diagnostics outside timing. The observation is therefore a pipeline-scope
comparison over five tested sizes, not a like-for-like end-to-end speedup or
evidence of asymptotic complexity. Python sort has the lowest median call time
at every tested size. Full numerical results and variability information are
reported in \cref{chap:results}, with interpretation limits collected in
\cref{chap:limitations}.

\section{Thesis Structure}
\label{sec:introduction-structure}

\Cref{chap:background} defines Jordan sorting, explains the nested family
structure, and places the 1986 and 1990 algorithms in context.
\Cref{chap:algorithm} develops the implementation-facing reconstruction of the
incremental procedure. \Cref{chap:implementation} describes the concrete
partial-order, sibling-list, audit, and execution-policy components.
\Cref{chap:methodology} specifies the frozen valid-input experiment, timing
boundaries, statistical units, and evidence contract. \Cref{chap:results}
presents correctness, runtime, variability, replication, and exploratory
findings. \Cref{chap:limitations} states the backend, timing, sampling, and
inferential boundaries. Finally, \cref{chap:conclusion} summarizes the answers
to the research questions and identifies work that remains outside the thesis.
